\documentclass[../../main.tex]{subfiles}
\begin{document}

{\bf [解答]}

\paragraph{(1)}

$\alpha = 100P + 10Q + R$ ($P, Q, R$ は $0 \sim 9$ までの整数) とおける．
この時題意より$\alpha' = 100R + 10Q + P$ だから
\begin{align}
  \alpha - \alpha'
  = 100(P - R) + (R - P)
  = 99(P - R) \label{eq:1}
\end{align}
である．$P, R$ は $0$ 以上 $9$ 以下の整数だから，$P - R$ は $-9$ 以上 $9$ 以下の全ての整数をとるので，$A$ のとり得る値の個数は $9 - (-9) + 1 = 19$個である．このうち正のものは $P - R$ が $1 \sim 9$ のものだから，$B$ は $9$個である．以上まとめて
\begin{align}
  A = 19 \text{個}, \quad
  B = 9 \text{個}
\end{align}
である．

\paragraph{(2)}

\cref{eq:1}から総和 $S$ として
\begin{align}
  S = \sum_{k=1}^9 99k = \dfrac{99}{2} \cdot 9 \cdot 10 = 4455
\end{align}
である．

\newpage

\end{document}
